\section{Introduction}

\begin{figure*}[h]
    \centering
    \includegraphics[width=1\textwidth]{figures/concept_fig.pdf}
    \caption{The two steps of the \method framework: \mine (in-the-wild user-written jailbreak tactics) and \compose (jailbreak tactics into diverse adversarial attacks).}
    \label{fig:concept_fig}
\end{figure*}

Despite ongoing efforts to enhance their safety, frontier LLMs remain vulnerable against unsafe user queries, especially adversarial attacks \citep{adversariallyaligned, zou2023universal}. The fact that models can be easily jailbroken raises significant concerns among researchers and policymakers ~\citep{hendrycks2023overview, biden2023executive,anwar2024safetychallenges}, motivating the research for systematically discovering and guarding against potential jailbreaks. In this work, we introduce the \method framework to address two challenges: 1) broadly identifying jailbroken behaviors of LLMs and 2) creating a publicly open, large-scale safety training resource for systematic defense. This resource is designed to help models robustly guard against \textit{vanilla} and \textit{adversarial} harmful user queries without causing over-refusal of benign queries or diminishing model general capabilities.

\textbf{The first challenge that \method addresses is to reveal vulnerabilities of LLMs against adversarial jailbreaks with scale and diversity.} We introduce \method, a practical red-teaming framework that composes automatically mined human-devised jailbreak tactics to transform vanilla harmful queries into many varieties of challenging adversarial attacks. \method improves over previous methods by diversifying the range of successful attack candidates while maintaining low computational costs, making it practical for scaling up. \method uncovers model vulnerabilities through a two-stage process: \textit{mining jailbreak tactics from in-the-wild} (\itwshort) \textit{chatbot logs} (\mine) and \textit{composing mined tactics into diverse adversarial attacks} (\compose). 

In the \mine stage, \method automatically maps out previously under-explored spaces of jailbreak tactics, significantly expanding the current taxonomy. To do so, it identifies 105K human-devised jailbreak tactics (5.7K unique clusters) from real-world user-chatbot interactions in \lmsys \citep{zheng2023lmsyschat1m} and \wildchat \citep{zhao2024inthewildchat}. In the \compose stage, \method generates diverse adversarial attack candidates by combining different selections of tactics using off-the-shelf LLMs like \mixtral \citep{jiang2024mixtral} and \gptfour \citep{openai2024gpt4}. It further refines attacks through lightweight off-topic and low-risk pruning to enhance attack quality and efficiency. With a suite of newly defined \textit{diversity} evaluation metrics, \method identifies up to 4.6 times more unique successful attacks against black-box and white-box LMs in 40\% fewer attack attempts compared to other state-of-the-art jailbreak methods, which sometimes struggle to find even two unique successful attacks.

\textbf{The second challenge \method addresses is to enhance open resources for safety training.} We apply \method to create \dataset, a large-scale, high-quality synthetic safety instruction-tuning data resource with 262K prompt and response pairs. \dataset contains four \textit{contrastive} components: 1) \textbf{vanilla harmful} queries conveying explicit unsafe requests across widespread risk categories, \eg malicious uses, harmful language \citep{weidinger2022taxonomy}; 2) \textbf{vanilla benign} queries that are similar to unsafe queries in form but convey no harmful intent, used to mitigate models' exaggerated safety behaviors \citep{safetytuned_llama}; 3) \textbf{adversarial harmful} queries that are jailbreaking versions of vanilla harmful queries converted by the \method heuristic; 4) \textbf{adversarial benign} queries used to counteract adversarial exaggerated safety behaviors, also generated by \method. \dataset is the first safety training resource to simultaneously address all four components, significantly improving upon existing resources with both enhanced scale and quality \citep{hh-rlhf,hh-rlhf-preference,safetytuned_llama,safe-rlhf}.


The unique composition and size of \dataset allow us to conduct extensive safety training experiments to study the scaling effect of safety training data and the interplay of data properties and model capabilities. Our experiments confirm the necessity of all components of \dataset for achieving \textit{balanced} safety behaviors, \ie robust safeguard without over-refusal on both vanilla and adversarial cases. Moreover, by mixing varying sizes of \dataset with \tulumix \citep{ivison2023tulu}, an instruction-tuning resource for teaching models general instruction-following and reasoning capabilities, we show that larger sizes of safety training data lead to gradually improving vanilla and adversarial safety features without sacrificing models' general capabilities, even at a scale orders of magnitude larger than studied in previous literature, measured by 15+ downstream tasks. Finally, although training on either vanilla or adversarial data improves performance on the other data type, the most robust safeguard comes with the hybridization of both. Our safety training insights pave the way towards building more transparent and safer future models.
